\section{比較と考察}
\label{sec:compare}

\begin{figure}[H]
\centering
\includegraphics[width=0.92\linewidth]{compare/M0-01_overview_map.pdf}
\caption[広域図（1 月）]{1 月後半、\SI{200}{\hecto\pascal} の平均風と 3 つの対象領域。
矢印の長さは 3 領域で共通の縮尺である。
ジェットの核が和歌山・愛媛の緯度帯にあり、北海道側では弱まっているのが見える。
和歌山と愛媛の領域は \ang{134}--\ang{136}E で重なっている。}
\label{fig:m0-jan}
\end{figure}

\begin{figure}[H]
\centering
\includegraphics[width=0.92\linewidth]{compare/M0-07_overview_map.pdf}
\caption[広域図（7 月）]{同じく 7 月後半。冬とは逆に、北海道側のほうが風が強い。
和歌山・愛媛の上空は風が弱く、矢印は南から南東を向く（北寄りの風）。}
\label{fig:m0-jul}
\end{figure}

\begin{figure}[H]
\centering
\includegraphics[width=\linewidth]{compare/F11_region_comparison.pdf}
\caption[北海道と和歌山の比較]{上 2 段は同じカラースケールでの平均風速、
下段はその差（北海道 $-$ 和歌山）。
赤は北海道のほうが強いビン・高度、青は和歌山のほうが強いビン・高度である。}
\label{fig:f11}
\end{figure}

\begin{figure}[H]
\centering
\includegraphics[width=\linewidth]{compare/F11b_region_comparison.pdf}
\caption[和歌山と愛媛の比較]{同じ図を和歌山と愛媛で描いたもの。
上 2 段はほとんど見分けがつかない。下段の差（和歌山 $-$ 愛媛）は
最大でも \SI{1.7}{\metre\per\second} で、カラースケールの幅が
図\ref{fig:f11} の 20 分の 1 以下であることに注意。}
\label{fig:f11b}
\end{figure}

\subsection{ジェットは季節とともに南北に動く}

図 \ref{fig:f11} の差分パネルに、両者の関係がそのまま出ている。

\begin{itemize}
  \item \textbf{冬（11--3 月）}、高度 \SIrange{8}{16}{\kilo\metre} は濃い青、
        すなわち和歌山のほうが最大 \SI{40}{\metre\per\second} 強い（2 月前半、\SI{200}{\hecto\pascal}）。
  \item \textbf{夏から初秋（7--9 月）}、同じ高度帯が赤に変わり、
        北海道のほうが最大 \SI{21}{\metre\per\second} 強い（9 月後半、\SI{200}{\hecto\pascal}）。
\end{itemize}

これは東アジアジェットの季節移動そのものである。
冬にはジェットの核が \ang{30}--\ang{35}N まで南下して和歌山と愛媛の真上を通り、
夏には \ang{40}N 以北へ退いて北海道の上に来る。

ここから読み取れるのは、\textbf{両者が逆位相で動いている}ということである。
一方が荒れているとき、もう一方は穏やかになる。
年間で最も穏やかな時期は、北海道が 7 月後半、和歌山と愛媛が 8 月後半と、
いずれも夏である。
これは対流圏のジェットが弱まる時期が共通しているためで、
差が出るのはその前後の季節である。

一方、\textbf{和歌山と愛媛はほとんど同じ気候である}（図\ref{fig:f11b}）。
経度で \ang{2} ほど離れていても、同じ緯度帯のジェットの下にあるので、
平均風速の差はどのビン・高度でも \SI{1.7}{\metre\per\second} 以内に収まる。
3 領域の違いを決めているのは経度ではなく緯度である。

\begin{table}[H]
\centering
\caption{主な数値の比較（領域平均、2016--2025 年）。
方位は機体が進む向きで、季節の値は正味の変位から求めた。}
\label{tab:compare}
\small
\begin{tabular}{lrrr}
\toprule
 & 北海道 & 和歌山 & 愛媛 \\
\midrule
格子点数 & 570 & 72 & 108 \\
\SI{200}{\hecto\pascal} 冬の平均風速 [\si{\metre\per\second}] & 40.2 & \textbf{75.5} & \textbf{75.4} \\
\SI{200}{\hecto\pascal} 夏の平均風速 [\si{\metre\per\second}] & \textbf{29.9} & 22.7 & 21.9 \\
\SI{200}{\hecto\pascal} 秋の平均風速 [\si{\metre\per\second}] & \textbf{44.5} & 40.3 & 39.4 \\
\SI{10}{\hecto\pascal} 夏の平均風速 [\si{\metre\per\second}] & 13.6 & \textbf{18.2} & \textbf{18.5} \\
領域内の風速差の最大 [\si{\metre\per\second}] & \textbf{31.7} & 12.0 & 13.3 \\
飛行距離 年間平均 [\si{\kilo\metre}] & 143 & 158 & 156 \\
飛行距離 最短のビン [\si{\kilo\metre}] & 62（7 月後半） & 49（8 月後半） & 50（8 月後半） \\
飛行距離 最長のビン [\si{\kilo\metre}] & 215（12 月前半） & 276（2 月前半） & 275（2 月前半） \\
冬（12--2 月）の飛行方位 & 東（\ang{83}） & 東（\ang{87}） & 東（\ang{86}） \\
7 月後半--8 月の飛行方位 & 東〜東南東（\ang{80}--\ang{123}） & \textbf{西寄り（\ang{247}--\ang{285}）} & \textbf{西寄り（\ang{252}--\ang{281}）} \\
ISA からのずれ 最大 [\si{\metre}] & \textbf{1036} & 787 & 790 \\
\bottomrule
\end{tabular}
\end{table}

\subsection{夏の飛行方位が逆を向く}

回収計画を左右するのは、\textbf{夏の飛行方位}である。

北海道は夏でも東から東南東へ飛ぶ（7 月後半--8 月で \ang{80}--\ang{123}）。
弱いながら対流圏上部に西風が残るためである。
和歌山と愛媛は、7 月後半から 8 月に西寄り（\ang{247}--\ang{285}）へ飛ぶ。
対流圏上部の風が弱まり、成層圏の東風（\SI{10}{\hecto\pascal} で
\SI{18}{\metre\per\second}、定常度 0.99）と釣り合うか、わずかに上回るからである。
7 月前半は東、9 月前半は北東と、半月ごとに向きが入れ替わるのもこの釣り合いのためである。

そのため南の 2 領域では、夏の\textbf{平均落下点は打ち上げ地点のすぐ近くに来る}。
1 回ごとには \SIrange{49}{58}{\kilo\metre} 飛ぶのに、東向きの日と西向きの日が打ち消し合う。
愛南町から上げると、7 月前半の平均落下点は陸、7 月後半は豊後水道の海上になる
（表\ref{tab:ainan-offshore}）。
海上への流出が主な懸念である北海道とは、
回収計画も安全評価も設計が変わる。

\subsection{領域を代表点で語れるかどうかも違う}

領域内の風速差は北海道で最大 \SI{31.7}{\metre\per\second}、
和歌山で \SI{12.0}{\metre\per\second}、愛媛で \SI{13.3}{\metre\per\second}。
和歌山と愛媛は緯度幅が \ang{2.5}--\ang{3} と狭くジェットの南北勾配をまたがないので、1 点で近似してよい。
北海道は \ang{6.5} の緯度幅があり、冬のジェット勾配をまたぐので、
\textbf{道南と道北を同じ数字で語ってはいけない}。

これは「解像度を上げたら何が分かったか」の具体例でもある。
\ang{1.25} 格子では北海道の領域に 54 点しか入らず、
この南北勾配を読み取るのは難しかった。

\subsection{上層の定常度はどの領域でも低い}

破裂高度帯（\SI{10}{\hecto\pascal}、約 \SI{31}{\kilo\metre}）の定常度は、
夏こそ 3 領域とも 0.99 だが、冬と春は 0.36--0.44、秋は 0.31--0.52 に落ちる。
ビン単位では 0.2 を切るところがあり、
北海道は 10 月前半（0.15）と 5 月前半（0.17）、
和歌山と愛媛は 3 月後半（どちらも 0.05）が最も低い。

この時期は\textbf{平均風向に意味がない}ので、
落下位置の平均値も参考程度にとどめるべきである。
散らばりのほうを見ること（本統計では p90 と楕円がそれにあたる）。

\subsection{Open-Meteo の上限による誤差はどれくらいか}
\label{sec:ceiling}

第 \ref{sec:intro-ceiling} 節で触れた「\SI{30}{\hecto\pascal} より上が分からない」という問題は、
本報告の統計を使えば定量化できる。
飛行モデルの層滞在時間（式 \eqref{eq:drift} の $T_k$）で重み付けし、
\SI{23.8}{\kilo\metre} より上の層が水平移動にどれだけ寄与するかを求めた。

\begin{table}[H]
\centering
\caption{\SI{23.8}{\kilo\metre}（Open-Meteo の上限）より上の層が、
水平移動に占める割合。1 つのビンの値ではなく、各月 2 ビンの平均。}
\begin{tabular}{lrrrrrr}
\toprule
 & 1 月 & 4 月 & 6 月 & 7 月 & 8 月 & 11 月 \\
\midrule
北海道 & 4.6 \% & 4.8 \% & 18.4 \% & \textbf{53.5 \%} & 20.5 \% & 7.1 \% \\
和歌山 & 1.1 \% & 3.6 \% & 20.0 \% & \textbf{111 \%} & 101 \% & 1.9 \% \\
愛媛   & 1.0 \% & 3.4 \% & 21.7 \% & \textbf{140 \%} & 86 \% & 2.0 \% \\
\bottomrule
\end{tabular}
\end{table}

冬は 1--5 \% にとどまる。
対流圏のジェットが桁違いに強いので、その上の層が加わっても正味の変位はほとんど変わらない。
\textbf{冬に関しては、既存実装の上限は実害がほとんど無い}。

一方、夏は状況が違う。
和歌山の 7 月に 111 \%、愛媛の 7 月に 140 \% という値が出ているのは、
分母が\textbf{正味の}変位で、分子が上層だけの寄与だからである。
上層の東風が下層の流れを打ち消しているため、正味の変位が小さくなり、比が 1 を超える。
言い換えると、\textbf{夏の和歌山と愛媛では、上層を無視すると向きそのものが逆に出る}。

そして第 \ref{sec:hokkaido} 章から第 \ref{sec:ehime} 章で見たとおり、
飛行距離が年間で最も短いのは 3 領域とも夏である。
打ち上げに最も向いた季節に、既存実装の誤差が最も大きくなる。

\subsection{モンテカルロの σ をどこから取るか}

本統計の副産物として、
シミュレータのモンテカルロ計算に与える風の摂動を実測から決められるようになった。

\begin{itemize}
  \item 高度ごとの風のばらつきは \texttt{u\_var}, \texttt{v\_var}, \texttt{uv\_cov}
        に入っている。等方ガウスではなく、実際の異方性を持った楕円が使える。
  \item 落下位置そのもののばらつきは \texttt{drift\_east\_var} などにある。
        こちらは各解析時刻ごとに飛行を積分した結果なので、
        層間の相関を仮定せずに得られている。
\end{itemize}

現在の実装（\texttt{frontend/src/mc/sampling.ts}）は全層に等方の摂動を与えているが、
実際のばらつきは高度によって桁で違い、向きにも偏りがある。
そこを置き換える材料はこの出力に揃っている。
